\section{Introduction}
The planar Mumford--Shah problem balances smoothness of a function against
the length of a set across which it may jump.  Its regularity conjecture
predicts that, in the interior, this set consists of regular arcs with only
free endpoints and triple junctions.  We prove this interior conclusion
for the absolute comparison problem defined below.  The proof also gives
local finiteness of the global connected components of the closed set.

\subsection{The variational problem}
\label{sec:problem}
Let $\Omega\subset\mathbb R^2$ be a bounded Lipschitz domain and let
$g\in L^\infty(\Omega)$ be real-valued.  We write $\mathcal H^1$ for
one-dimensional Hausdorff measure.  A pair $(u,K)$ is \emph{admissible} if
\[
 K\subset\Omega\text{ is relatively closed},\qquad
 \mathcal H^1(K)<\infty,\qquad
 u\in L^2(\Omega),\qquad
 u|_{\Omega\setminus K}\in W^{1,2}(\Omega\setminus K).
\]
No rectifiability of $K$ is assumed.  Values of $u$ on $K$ do not affect
any integral.  For open $V\subset\Omega$, set
\begin{equation}\label{eq:original-energy}
 \mathcal F_g(u,K;V)
 =\int_{V\setminus K}|\nabla u|^2\,dx
   +\mathcal H^1(K\cap V)+\int_V|u-g|^2\,dx.
\end{equation}
Here and below, $V\Subset\Omega$ means that $\overline V$ is a compact
subset of $\Omega$.

An admissible pair is an \emph{absolute minimizer} if, for every open
$V\Subset\Omega$ and every admissible $(w,L)$ with
\begin{equation}\label{eq:compact-comparison}
 (L\mathbin{\triangle}K)\cup\operatorname{spt}(w-u)\Subset V,
\end{equation}
one has $\mathcal F_g(u,K;V)\le\mathcal F_g(w,L;V)$.
The support of a function is understood in the essential sense.
The pair is \emph{reduced} if
\begin{equation}\label{eq:reduced}
 K=\operatorname{spt}_{\Omega}(\mathcal H^1\!\llcorner K).
\end{equation}
Equivalently, every disk centered at a point of $K$ and compactly
contained in $\Omega$ contains a positive length of $K$.  This removes
closed additions of zero length that have no effect on the energy.

\subsection{Main theorem}
\begin{theorem}\label{thm:main}
Let $(u,K)$ be a reduced absolute minimizer of
\eqref{eq:original-energy} in a bounded Lipschitz domain
$\Omega\subset\mathbb R^2$, with $g\in L^\infty(\Omega)$.
Every $x\in K$ has an interior neighborhood in which $K$ has one of the
following forms, with $C^{1,\alpha}$ arcs for some $\alpha>0$:
\begin{enumerate}[label=\textup{(\roman*)},nosep]
 \item a single embedded arc with $x$ in its interior;
 \item a single embedded arc ending at $x$;
 \item three embedded arcs meeting only at $x$, with pairwise angles
       $120$ degrees.
\end{enumerate}
Moreover, for every open $U\Subset\Omega$,
\begin{equation}\label{eq:component-finiteness}
 \#\{C\in\pi_0(K):C\cap U\ne\varnothing\}<\infty,
\end{equation}
where $\pi_0(K)$ denotes the set of global connected components of $K$.
\end{theorem}

The theorem gives the interior regularity conclusion of the planar
Mumford--Shah conjecture.  The compact-set qualification is essential to
the scope of the statement: we make no assertion about finiteness up to
$\partial\Omega$.  The components counted in
\eqref{eq:component-finiteness} are components of $K$, rather than of a
truncation $K\cap U$ or of its complement.

The geometric conclusion also gives an endpoint integrability estimate for
the gradient.  Corollary~\ref{cor:weak-l4} proves
\[
 \nabla u\in L^{4,\infty}_{\mathrm{loc}}(\Omega),
 \qquad \nabla u\in L^p_{\mathrm{loc}}(\Omega)\quad(0<p<4).
\]
Here weak $L^4$ means that the area of the gradient superlevel set
$\{|\nabla u|>t\}$ in each relatively compact region is bounded by a
constant times $t^{-4}$.  The exponent reflects the inverse-square-root
gradient at a crack tip.

\subsection{History and related work}
Mumford and Shah introduced their functional as a model for image
segmentation \cite{mumford_shah1989}. The fidelity term keeps the
reconstruction close to the observed image, the Dirichlet term favors
smooth variation within regions, and the length term penalizes the edges
between them. Their regularity conjecture asks whether minimization itself
forces the simple interior geometry suggested by this model. Deriving that
geometry requires controlling arbitrary finite-length discontinuity sets,
including possible accumulation of their components.

The theory of special functions of bounded variation provides the weak
variational framework for this problem. De Giorgi and Ambrosio introduced
this free-discontinuity setting \cite{degiorgi_ambrosio1988}, and Ambrosio
developed the compactness and existence theory
\cite{ambrosio1989,ambrosio1990}. The existence and essential-closure results
of De Giorgi, Carriero, and Leaci, together with the Dirichlet theory of
Carriero and Leaci, connect weak minimizers with the classical closed-set
formulation \cite{degiorgi_carriero_leaci1989,carriero_leaci1990}.
We use this connection to establish rectifiability of the given set before
applying geometric regularity theory. The underlying bounded-variation
extension and trace calculus are treated systematically in
Ambrosio, Fusco, and Pallara \cite{ambrosio_fusco_pallara2000}.

Bonnet introduced the planar blow-up and monotonicity approach and
classified global minimizers with connected crack set
\cite[Theorems~2.2, 3.1, and~4.1]{bonnet1996}.
Regular-arc and partial-regularity theory was developed independently by
David and by Ambrosio, Fusco, and Pallara
\cite{david1996,ambrosio_fusco_pallara1997}; David also established the
triple-junction theory and developed its quantitative geometric framework
\cite{david1996,david2005}. Bonnet and David proved global minimality of
the crack tip and classification when the part outside one connected
component is bounded \cite{bonnet_david2001}. David and L\'eger proved
that a global minimizer whose crack disconnects the plane must be a line
or a $120$-degree triod \cite[Corollary~9.16]{david_leger2002}.
In the generalized-limit formulation recalled below, these classification
results leave the possible bounded components of a limiting crack as the
remaining obstruction. Higher-integrability theorems of De Lellis--Focardi
in the plane and De Philippis--Figalli in arbitrary dimension supplied
further control of the exceptional set
\cite{delellis_focardi2013,dephilippis_figalli2014}.

Regularity at an endpoint requires additional analysis. Work of Andersson
and Mikayelyan on crack tips \cite{andersson_mikayelyan2019} was revisited
by De Lellis, Focardi, and Ghinassi
\cite{delellis_focardi_ghinassi2021,delellis_focardi_ghinassi_erratum}.
For the bounded measurable fidelity datum used here, we invoke the exact
geometric regularity and compactness statements in the monograph of
De Lellis and Focardi \cite{delellis_focardi2025}. Recent work of Labourie
and Lemenant establishes at most three limiting values and controls local
components of the complement \cite{labourie_lemenant2026}. The component
count in Theorem~\ref{thm:main} concerns the closed crack set itself.

Deangelis's recent preprint gives a proof of the global classification and
the resulting interior arc, endpoint, and triple-junction structure
\cite[Theorems~1.1 and~3.4]{deangelis2026}. His compact-translation
argument combines relaxed second variations, equality in a planar Hodge
identity, and holomorphic rigidity. Both arguments use whole-plane
projections and value variations with independently prescribed one-sided
traces on a regular arc; see \cite[Section~3.2]{deangelis2026}.
The proof below obtains a scalar harmonic interaction from finite
translations and an explicit cutoff at infinity; minimality forces this
interaction to be constant.

\subsection{Proof strategy}
We first pass from the stated finite-length formulation to a locally
bounded, rectifiable minimizing pair.  The reduction uses clipping inside
disks and a Dirichlet comparison with smooth approximations of the boundary
trace; it preserves the original pair.

At an interior point $x$, the rescaled closed sets $(K-x)/r_j$, with
$r_j\downarrow0$, have global generalized limits $H$.  These are limits
of absolute minimizers, with additive constants in each complementary
component recorded as part of the limiting data.  The established
compactness and rigidity theorems reduce their classification to excluding
a bounded component when $\mathbb R^2\setminus H$ is connected.  In this
case the limit function $v$ is an absolute minimizer of the zero-fidelity
energy on the whole plane.

A bounded component can be enlarged to a compact piece $A\subset H$ at
positive distance from $B=H\setminus A$.  We split the gradient of $v$
into a field carrying the jumps across $A$ and the gradient of a harmonic
potential that extends through $A$.  Translating $A$ leaves its length unchanged.  The
interaction between these two fields is harmonic in the translation
parameter; minimality forces it to be constant.  Independent changes in
the jump values along a regular subarc of $A$ then force this harmonic
potential to be affine.  The energy growth bound makes its gradient zero;
the value-variation identity then makes the remaining field zero.
Removing $A$ saves positive length, a contradiction.

The known classification now gives only a line, a halfline or three rays
for a nonempty blow-up set.  Epsilon regularity transfers these models to
the original minimizer.  Finally, a finite cover of $K\cap\overline U$ by
connected model neighborhoods proves \eqref{eq:component-finiteness}.
The resulting absence of irregular points also yields the local weak-$L^4$
gradient estimate through the De Lellis--Focardi equivalence.

The reduction occupies Section~\ref{sec:initial-reduction}, and
Section~\ref{sec:known-facts} records the compactness and regularity
results used below.  Sections~\ref{sec:isolation} and~\ref{sec:compact-piece} establish
the compact-piece exclusion.  Section~\ref{sec:assembly} completes the
classification, proves Theorem~\ref{thm:main}, and derives the weak-$L^4$
consequence.
